\documentclass[11pt]{article}

\usepackage[margin=1in]{geometry}
\usepackage{amsmath}
\usepackage{amssymb}
\usepackage{booktabs}
\usepackage[hidelinks]{hyperref}

\title{Revised Statistical Calibration for Evaluator-Based Scoring}
\author{Statistical Evals Project}
\date{October 21, 2026}

\begin{document}

\maketitle

\begin{abstract}
The initial statistical calibration method successfully measured evaluator
noise and produced more stable composite scores. However, the Tau3 experiment
showed that statistical reliability is not the same as semantic relevance.
The learned score was more repeatable than an equal-weight score, but it was
less aligned with Tau3's native task-success reward. This document proposes a
revised framework that separates measurement uncertainty from evaluator
importance, learns target-specific weights when labels are available, and
treats customer-declared priorities as the primary definition of a
multi-dimensional quality score. The revised framework may also substantially
reduce the number of agent executions and evaluator calls required from
customers.
\end{abstract}

\section{Background}

An evaluation report contains scores from multiple evaluators. These
evaluators may differ in:

\begin{itemize}
    \item judge replay noise;
    \item score granularity;
    \item sensitivity to different cases and agent outputs; and
    \item semantic relevance to the customer's actual objective.
\end{itemize}

The initial project attempted to replace an equal-weight average with a
reliability-weighted composite. For evaluator \(e\), the provisional
reliability was:

\begin{equation}
r_e =
\frac{v_{c,e}}
     {v_{c,e} + v_{a,e} + v_{j,e}/k},
\end{equation}

where:

\begin{itemize}
    \item \(v_{c,e}\) is variance across cases;
    \item \(v_{a,e}\) is variance across agent runs of the same case;
    \item \(v_{j,e}\) is variance across judge replays of a fixed run; and
    \item \(k\) is the number of judge replays.
\end{itemize}

The active weighting rule was:

\begin{equation}
w_e =
\frac{r_e / \sqrt{\mathrm{step}_e}}
     {\sum_t r_t / \sqrt{\mathrm{step}_t}}.
\end{equation}

This rule was intended to reward repeatable evaluators while preventing
coarse evaluators from dominating the composite.

\section{Tau3 Experiment}

\subsection{Collection}

The main Tau3 experiment contained:

\begin{itemize}
    \item 80 cases across four domains;
    \item 3 independent agent runs per case;
    \item 3 judge replays per fixed output;
    \item 3 evaluators; and
    \item 2,160 total evaluator observations.
\end{itemize}

The evaluator set was:

\begin{itemize}
    \item GoalSuccessRateEvaluator;
    \item HelpfulnessEvaluator; and
    \item InstructionFollowingEvaluator.
\end{itemize}

Tau3's binary native reward was excluded from the weighting calculation. It
was retained as an external task-success signal for evaluating whether the
resulting composite remained aligned with objective task completion.

\subsection{Calibration Size}

The analysis repeatedly subsampled cases, agent runs, and judge replays. The
smallest tested design that consistently passed the provisional uncertainty
and stability thresholds was:

\begin{center}
\[
80\ \text{cases}
\times 3\ \text{agent runs}
\times 2\ \text{judge replays}.
\]
\end{center}

Designs with 12 or 20 cases were too uncertain. Designs with 40 cases
occasionally passed, but not reliably. Adding a third judge replay at 80 cases
narrowed the weight intervals but did not materially improve the uncertainty
of the final composite score.

\subsection{Learned Weights}

The complete fit produced approximately the following results:

\begin{center}
\begin{tabular}{lrr}
\toprule
Evaluator & Reliability & Learned weight \\
\midrule
Goal success & 0.699 & 0.238 \\
Helpfulness & 0.711 & 0.612 \\
Instruction following & 0.482 & 0.150 \\
\bottomrule
\end{tabular}
\end{center}

Helpfulness received most of the weight because it had high estimated
reliability and a smaller declared step size. The inverse-square-root step
adjustment amplified its contribution relative to the two binary evaluators.

\subsection{Stability and Held-Out Reward}

Learned weighting reduced score variation across replay subsamples:

\begin{center}
\begin{tabular}{lr}
\toprule
Method & Score standard deviation \\
\midrule
Equal weighting & 0.0036 \\
Learned weighting & 0.0028 \\
\bottomrule
\end{tabular}
\end{center}

This is approximately a 24 percent reduction in score variation. However, the
held-out association with Tau3 task success moved in the opposite direction:

\begin{center}
\begin{tabular}{lr}
\toprule
Method & Correlation with native reward \\
\midrule
Equal weighting & 0.139 \\
Learned weighting & 0.107 \\
\bottomrule
\end{tabular}
\end{center}

The learned composite was more stable, but it was a weaker proxy for Tau3
task completion. Both correlations were weak.

\subsection{Shrinkage Experiment}

The follow-up experiment blended equal and learned weights:

\begin{equation}
w_e^{(\alpha)}
=
(1-\alpha)w_e^{\mathrm{equal}}
+ \alpha w_e^{\mathrm{learned}},
\qquad 0 \leq \alpha \leq 1.
\end{equation}

\begin{center}
\begin{tabular}{rrr}
\toprule
Learned fraction \(\alpha\) & Score standard deviation
& Reward correlation \\
\midrule
0.00 & 0.0036 & 0.139 \\
0.25 & 0.0033 & 0.133 \\
0.50 & 0.0030 & 0.126 \\
0.75 & 0.0028 & 0.117 \\
1.00 & 0.0028 & 0.107 \\
\bottomrule
\end{tabular}
\end{center}

Increasing the learned fraction continuously improved stability while
continuously reducing task-reward alignment. No tested blend improved both
measurements.

\section{What the Experiment Revealed}

\subsection{Reliability Is Not Importance}

Reliability describes whether an evaluator is repeatable and discriminative.
It does not determine whether that evaluator measures the quality dimension a
customer cares about.

An evaluator can be:

\begin{itemize}
    \item highly consistent but irrelevant to the target;
    \item noisy but semantically essential;
    \item fine-grained but off-axis; or
    \item coarse but closely aligned with task success.
\end{itemize}

The current weighting rule treats repeatability and score resolution as
importance. The Tau3 result shows that this assumption is not generally
defensible.

\subsection{Agent-Run Variation Is Not Pure Evaluator Noise}

Judge replay variation is evaluator measurement noise because the underlying
agent output is fixed. Agent-run variation is different: a new run produces a
new output and may represent real agent behavior.

Penalizing an evaluator for detecting variation between agent runs can
incorrectly suppress useful signal. Agent-run variation should primarily
contribute to uncertainty about the agent's expected performance, rather than
automatically reducing evaluator importance.

\subsection{Step Size Is Not Semantic Priority}

A fine-grained scale offers more resolution, but resolution is not the same as
importance. Dividing directly by the square root of step size can give a
fine-grained evaluator substantial influence even when it measures a secondary
quality dimension.

Step size should instead be represented through the evaluator's measurement
model or uncertainty. It should not independently determine semantic weight.

\section{Different Objectives Produce Different Weights}

Before selecting a weighting formula, the system must define what the final
number is supposed to accomplish. There is no universally optimal weight
vector because different objectives answer different questions.

\subsection{Objective 1: Measure One Construct Precisely}

Suppose several evaluators are intended to measure the same underlying
quantity. For example, three independent judges may all estimate task
correctness. In this setting, the statistical objective may be to minimize the
variance of their combined estimate.

If the evaluators are unbiased measurements of the same construct and their
measurement-error covariance matrix is \(\Sigma\), the unconstrained
minimum-variance weights are proportional to:

\begin{equation}
w^{\mathrm{precision}}
=
\frac{\Sigma^{-1}\mathbf{1}}
     {\mathbf{1}^{\mathsf{T}}\Sigma^{-1}\mathbf{1}}.
\end{equation}

The intuition is that a noisy evaluator contributes less, while two evaluators
with strongly correlated errors provide less additional information than two
independent evaluators. The covariance matrix matters because merely adding
more judges does not help much when they make the same mistakes.

These weights are defensible only when the evaluators are genuinely measuring
the same construct. Applying them to Goal Success, Helpfulness, and
Instruction Following would silently assume that those dimensions are
interchangeable measurements of one latent variable. That is a semantic
assumption, not a consequence of the variance calculation.

\subsection{Objective 2: Predict an External Outcome}

If the desired score is meant to predict an external outcome, such as whether
the agent completes a task, then the objective is predictive accuracy.

The weights should answer:

\begin{quote}
Which combination of evaluator scores best predicts the outcome on cases that
were not used to fit the combination?
\end{quote}

An evaluator may receive a large predictive coefficient even if it is
moderately noisy, provided that it contains information about the target that
the other evaluators do not contain. Conversely, a perfectly repeatable
evaluator may receive little weight when it is unrelated to the target.

These are target-specific weights. Weights learned to predict task completion
must not be described as universal quality weights. A different target, such
as customer satisfaction or policy compliance, may produce a different
vector.

\subsection{Objective 3: Express Customer Priorities}

When evaluators measure different quality dimensions, the most direct and
defensible objective is often customer utility. The customer declares a
priority vector \(p\):

\begin{equation}
p_e \geq 0,
\qquad
\sum_e p_e = 1.
\end{equation}

The composite for observation \(n\) is then:

\begin{equation}
S_n^{\mathrm{declared}}
=
\sum_e p_e x_{ne}.
\end{equation}

For example, a customer might declare:

\begin{center}
\begin{tabular}{lr}
\toprule
Dimension & Declared weight \\
\midrule
Task success & 0.60 \\
Helpfulness & 0.25 \\
Instruction following & 0.15 \\
\bottomrule
\end{tabular}
\end{center}

These weights are not estimates of a hidden truth. They are an explicit
statement that one point of task success is worth more to this application
than one point of helpfulness or instruction following.

Statistics still play an important role. They determine how uncertain the
reported composite is, whether an evaluator is too noisy to support the
decision, and how much additional evidence is required. They do not silently
rewrite the customer's priorities.

\subsection{Objective 4: Make a Regression Decision}

A merge gate does not necessarily need an absolute universal score. Its
objective is to determine whether a candidate changed relative to a baseline.

For matched case \(i\) and evaluator \(e\), define the paired difference:

\begin{equation}
\Delta_{ie}
=
x_{ie}^{\mathrm{candidate}}
- x_{ie}^{\mathrm{baseline}}.
\end{equation}

The product can estimate a mean difference and confidence interval for each
quality dimension. It may then apply customer-declared rules such as:

\begin{itemize}
    \item task success must not regress;
    \item instruction following must remain above a minimum threshold; and
    \item helpfulness may trade off against latency within a declared range.
\end{itemize}

This is a decision problem rather than a measurement problem. A vector of
paired effects and constraints may be more defensible than forcing every
dimension into one learned scalar.

\subsection{Why the Objective Must Be Declared First}

The four objectives can legitimately produce different answers:

\begin{itemize}
    \item precision weighting asks which measurements are least noisy;
    \item predictive weighting asks which measurements forecast a target;
    \item declared weighting asks what the customer values; and
    \item comparison testing asks whether a candidate changed enough to affect
    a decision.
\end{itemize}

A weighting method is defensible only when its mathematical objective matches
the product claim attached to the resulting score.

\section{Revised Statistical Framework}

\subsection{Layer 1: Measurement}

The first layer estimates each evaluator's score and uncertainty without
deciding its importance.

Let:

\begin{equation}
s_{eiar}
=
\mu_e + c_{ei} + a_{eia} + \epsilon_{eiar},
\end{equation}

where:

\begin{itemize}
    \item \(e\) indexes evaluators;
    \item \(i\) indexes cases;
    \item \(a\) indexes agent runs;
    \item \(r\) indexes judge replays;
    \item \(c_{ei}\) is the evaluator-specific case effect;
    \item \(a_{eia}\) is the effect of a particular agent output; and
    \item \(\epsilon_{eiar}\) is judge replay noise.
\end{itemize}

This equation is a bookkeeping device for explaining where observed score
movement comes from. The term \(\mu_e\) is evaluator \(e\)'s typical score.
The case effect \(c_{ei}\) captures that some tasks are easier or harder. The
agent-run effect \(a_{eia}\) captures that the agent can produce a better or
worse output on another attempt. The replay term \(\epsilon_{eiar}\) is the
judge's wobble when the underlying output has not changed.

Only the final term is unambiguously evaluator measurement noise. The case and
agent-run terms describe real differences in what was evaluated.

For replay repeatability, both case and agent-output variation are signal:

\begin{equation}
\mathrm{repeatability}_e
=
\frac{v_{c,e} + v_{a,e}}
     {v_{c,e} + v_{a,e} + v_{j,e}/k}.
\end{equation}

This statistic describes whether repeated judging preserves distinctions
between fixed outputs. It should be reported as a diagnostic, not directly
treated as semantic importance.

The numerator asks how much fixed outputs genuinely differ from one another.
The denominator adds the movement caused by replaying the judge. A
repeatability near one means that most observed movement is attached to the
fixed outputs. A value near zero means judge wobble is large relative to the
differences the evaluator is trying to measure.

For a balanced experiment with \(N\) cases, \(A\) runs per case, and \(K\)
judge replays, the approximate uncertainty of an evaluator's suite mean is:

\begin{equation}
\mathrm{Var}(\widehat{\mu}_e)
\approx
\frac{v_{c,e}}{N}
+ \frac{v_{a,e}}{NA}
+ \frac{v_{j,e}}{NAK}.
\end{equation}

This decomposition determines whether more cases, agent runs, or judge
replays are needed. The denominators provide the intuition:

\begin{itemize}
    \item more cases average down uncertainty caused by task sampling;
    \item more agent runs average down uncertainty caused by stochastic agent
    behavior; and
    \item more judge replays average down judge measurement noise.
\end{itemize}

If the customer declares semantic weights \(p\), the uncertainty of the
composite is:

\begin{equation}
\mathrm{Var}
\left(
\sum_e p_e \widehat{\mu}_e
\right)
=
p^{\mathsf{T}}
\mathrm{Cov}(\widehat{\mu})
p.
\end{equation}

This equation is important because the same customer weight vector can be
retained while the confidence interval responds to evaluator noise and
correlation. Measurement statistics affect how certain the system is, not what
the customer values.

\subsection{Layer 2: Target-Specific Learning}

The second layer is optional. It learns a predictive score when the customer
has explicitly identified an external target. It does not replace
customer-declared weighting for a general multi-dimensional quality report.

If a benchmark reward, human label, or other external target \(y_n\) exists,
weights can be learned through a constrained supervised objective:

\begin{equation}
\min_{w}
\sum_n
\mathcal{L}
\left(
y_n,
\sum_e w_e x_{ne}
\right)
+ \lambda \lVert w-p \rVert_2^2
+ \gamma w^{\mathsf{T}}\Sigma_j w,
\end{equation}

subject to:

\begin{equation}
w_e \geq 0,
\qquad
\sum_e w_e = 1.
\end{equation}

Here:

\begin{itemize}
    \item \(x_{ne}\) is evaluator \(e\)'s aggregated score for observation
    \(n\);
    \item \(y_n\) is the external quality target;
    \item \(p\) is an equal-weight or customer-declared prior;
    \item \(\lambda\) controls shrinkage toward the prior; and
    \item \(\Sigma_j\) is the covariance matrix of judge measurement error.
\end{itemize}

The first term rewards combinations that predict the target on training data.
The second term prevents a small or noisy dataset from pulling the vector too
far from the declared prior. The third term penalizes combinations whose
weighted judge errors have high variance.

The covariance form is more general than penalizing each evaluator
independently. If judge errors are independent, it reduces to a sum of terms
like \(w_e^2 v_{j,e}/K\). If evaluators share a judge model or rubric failure,
the off-diagonal covariance terms represent their correlated errors.

The constraints also have practical meaning. Non-negative coefficients prevent
the model from treating a higher Helpfulness score as evidence of lower
quality. Requiring the weights to sum to one keeps the result on an
interpretable convex-combination scale.

Overall, the formulation separates three concerns:

\begin{itemize}
    \item prediction loss measures semantic relevance;
    \item the prior protects against unstable small-sample coefficients; and
    \item the noise penalty protects score repeatability.
\end{itemize}

\subsection{Customer-Declared Weights as the Default}

When dimensions represent distinct values rather than redundant measurements,
the recommended default is:

\begin{equation}
\widehat{S}
=
\sum_e p_e \widehat{\mu}_e,
\end{equation}

where \(p\) is supplied by the customer or application owner.

The system should help the customer understand the consequences of this
choice. It can report evaluator uncertainty, correlations between evaluators,
sensitivity of the composite to each weight, and whether one evaluator
dominates the result. It should not claim that replay statistics discovered
the customer's utility function.

Customers who do not want to specify priorities may use equal weights as a
transparent neutral convention. Equal weighting is still a policy choice, but
it is easier to explain than a data-derived vector whose objective is unclear.

\subsection{Binary Targets}

For a binary target such as Tau3 reward, a logistic model may be more
appropriate:

\begin{equation}
\Pr(y_n=1 \mid x_n)
=
\sigma
\left(
\beta_0 + \sum_e \beta_e x_{ne}
\right).
\end{equation}

The output is an estimated probability of task success rather than an
arbitrary weighted average.

Candidate methods should be selected using held-out proper scoring metrics,
including:

\begin{itemize}
    \item Brier score;
    \item log loss;
    \item probability calibration; and
    \item replay stability.
\end{itemize}

Correlation may remain a descriptive diagnostic, but it should not be the
primary objective for a binary outcome.

\subsection{When No External Target Exists}

Without a reward, human label, or declared objective, semantic importance is
not identifiable from evaluator replay data alone.

In this case, the system should:

\begin{itemize}
    \item ask the customer to declare semantic priorities when the dimensions
    differ;
    \item otherwise retain equal weights as a transparent fallback;
    \item report evaluator-specific uncertainty;
    \item warn about highly noisy evaluators;
    \item avoid presenting repeatability weights as correctness weights; and
    \item preserve the evaluator score vector alongside any composite.
\end{itemize}

\section{Multiple Quality Dimensions}

Task success, helpfulness, and instruction following may represent different
constructs. A single automatically learned scalar can obscure those
differences.

A report can instead preserve separate dimensions:

\begin{center}
\begin{tabular}{lr}
\toprule
Dimension & Example result \\
\midrule
Task success & \(0.82 \pm 0.04\) \\
Helpfulness & \(0.88 \pm 0.02\) \\
Instruction following & \(0.91 \pm 0.05\) \\
\bottomrule
\end{tabular}
\end{center}

Customers may then declare semantic priorities or gating constraints. For
example:

\begin{itemize}
    \item block if task success regresses;
    \item block if instruction following falls below 0.80; and
    \item use helpfulness as a secondary comparison.
\end{itemize}

Measurement reliability determines confidence in each dimension. Customer
priorities determine its importance.

\section{Revised Customer Workflow}

\begin{enumerate}
    \item Run the normal evaluation and store raw evaluator observations.
    \item Declare whether the product needs a multi-dimensional quality report,
    a target predictor, or a paired regression decision.
    \item For a quality composite, provide semantic weights or explicitly
    select equal weighting.
    \item Identify an external target when predictive calibration is desired,
    such as benchmark reward or human labels.
    \item Replay judges on a representative subset to estimate judge noise.
    \item Fit candidate objective models using case-grouped cross-validation.
    \item Compare held-out quality, uncertainty, and replay stability.
    \item Save a versioned profile only if the configured readiness criteria
    pass.
    \item Apply the compatible profile to current and historical reports.
    \item Recalibrate when the dataset, evaluators, judge model, or objective
    changes.
\end{enumerate}

The scoring states remain:

\begin{itemize}
    \item \texttt{uncalibrated}: no usable profile exists; use equal weights;
    \item \texttt{equal\_weight\_fallback}: calibration was attempted but the
    evidence was insufficient; and
    \item \texttt{calibrated}: a compatible profile passed its configured
    readiness criteria.
\end{itemize}

\section{Profile Updates and Order Invariance}

New evidence may arrive over multiple customer runs. The final profile should
not depend accidentally on the order in which equivalent observations arrived.

For a fixed observation set \(D\) and any permutation \(\pi(D)\), the desired
invariant is:

\begin{equation}
\mathrm{Fit}(D)
=
\mathrm{Fit}(\pi(D)).
\end{equation}

P0 should achieve this by storing the applicable observations and refitting the
profile from all of them. The implementation should use deterministic
observation ordering, deterministic seeds, and a deterministic optimizer.

An online update of the form:

\begin{equation}
w_{t+1}
=
\mathrm{Update}(w_t, B_{t+1})
\end{equation}

can be path-dependent because early batches influence the state inherited by
later batches. Fixed learning rates, early stopping, and confidence schedules
can all make two arrival orders produce different final vectors.

If a future implementation requires streaming updates, it should either use
mergeable sufficient statistics or periodically refit from the complete
applicable calibration window. Order dependence should be introduced only
deliberately, for example when recent data is intended to receive more weight
because the experiment is drifting. In that situation the time-decay rule must
be explicit and the resulting profile must be versioned.

\section{Potential Customer Cost Reduction}

\subsection{Current Experimental Cost}

The provisional Tau3 design required:

\begin{equation}
80\ \text{cases}
\times 3\ \text{agent runs}
\times 2\ \text{judge replays}
\times 3\ \text{evaluators}
= 1{,}440\ \text{evaluator calls}.
\end{equation}

It also required:

\begin{equation}
80 \times 3 = 240
\end{equation}

agent executions.

\subsection{Candidate Reduced-Cost Design}

Relevance learning and noise estimation do not require the same repeated
observations. A candidate design could use one ordinary evaluation over all
cases and replay judges only on a subset:

\begin{align}
\text{Base evaluation}
&=
80 \times 1 \times 1 \times 3
= 240\ \text{evaluator calls}, \\
\text{Noise audit}
&=
20 \times 1 \times 2 \times 3
= 120\ \text{additional calls}, \\
\text{Total}
&=
360\ \text{evaluator calls}.
\end{align}

This illustrative design uses 75 percent fewer evaluator calls and 67 percent
fewer agent executions than the provisional maximum. It has not yet been
validated as sufficient.

\subsection{Adaptive Collection}

The product should eventually collect additional evidence only when needed:

\begin{center}
\[
\text{initial evaluation}
\longrightarrow
\text{uncertainty check}
\longrightarrow
\text{targeted additional runs or replays}.
\]
\end{center}

Judge replays should concentrate on evaluators whose replay noise remains
uncertain. Additional agent runs should be used when agent stochasticity
dominates the uncertainty of the suite mean. Stable experiments should not
pay the maximum calibration cost repeatedly.

Once fitted, a compatible profile can be reused for historical and future
reports. Recalibration becomes periodic rather than a mandatory cost on every
evaluation.

\subsection{Cost Caveat}

The revised method shifts some cost toward obtaining a trustworthy target:

\begin{itemize}
    \item objective benchmark rewards are often inexpensive;
    \item historical customer outcomes may already exist;
    \item human quality labels may be more expensive than model replays; and
    \item without any external target, the system cannot learn semantic
    importance automatically.
\end{itemize}

The cost advantage is largest when customers already possess objective
verifier results or historical labels.

\section{Immediate Experiment Plan}

The next analysis can use the existing Tau3 collection without additional
model calls.

\begin{enumerate}
    \item Aggregate the judge replays for each of the 240 case-runs.
    \item Split training and validation folds by case, not by individual row.
    \item Compare:
    \begin{itemize}
        \item equal weighting;
        \item one or more customer-declared priority vectors;
        \item the existing \(r/\sqrt{\mathrm{step}}\) rule;
        \item precision weighting as a same-construct diagnostic;
        \item constrained supervised stacking; and
        \item supervised stacking with an equal-weight prior and judge-noise
        penalty.
    \end{itemize}
    \item Evaluate held-out Brier score, log loss, calibration, replay
    stability, and weight stability.
    \item Bootstrap complete cases to estimate uncertainty.
    \item Retain equal weighting if no alternative improves held-out quality
    reliably.
    \item Simulate smaller replay subsets to estimate the reduced-cost
    calibration budget.
\end{enumerate}

\section{Conclusion}

The Tau3 experiment did not invalidate statistical calibration. It clarified
what calibration can and cannot establish.

Variance decomposition is useful for measuring uncertainty and determining
how much evidence is needed. It cannot determine which quality dimensions
matter to a customer. The revised design therefore separates measurement
reliability from semantic importance.

For a general multi-dimensional quality score, customer-declared priorities
should determine the composite and statistical analysis should determine its
uncertainty. When an external predictive objective exists, the system may
learn target-specific weights under cross-validation while penalizing judge
noise and extreme departures from declared priorities. When neither a target
nor declared priorities exist, equal weighting should remain in place as a
transparent fallback.

\end{document}
